\chapter{Conclusion}
\label{sec:conclusion}

This thesis revisited operating system security through the lens of \emph{protection alignment}—the need to reconcile legacy OS abstractions with modern hardware-based isolation.
It argued that sustainable system security cannot rely on stronger mechanisms alone, but on the coherence of protection semantics across hardware and software layers.

Through two complementary systems, \capacity and \incognitos, this work demonstrated how such alignment can be achieved in practice.
\capacity retrofits capability-based protection into legacy process abstractions, using ARM Pointer Authentication (\gls{pa}) and Memory Tagging (\gls{mte}) to ensure consistent, hardware-backed enforcement within a process.
\incognitos redefines the OS’s role in trusted execution environments, treating it as an active mediator that mitigates side-channel leakage through scheduling and memory re-randomization.

Together, these systems show that aligning protection models enhances both the security and efficiency of computation, transforming hardware support from an external defense mechanism into an integral part of OS design.

Looking forward, this principle of protection alignment suggests a broader agenda for systems research: operating systems must evolve toward cooperative enforcement with hardware, bridging traditional abstractions with emerging isolation technologies to form a coherent, verifiable foundation for secure computation.
